% lit_capture.tex -- Literature 2.3: capture proof, docking and magnetic grasping (drafted 12 Oct 2026 from the
% thesis plan of 11 Oct). Moved here from the docking paragraph of control.tex and extended. Entries the plan lists
% that are not yet in the bibliography (Loianno 2018, Larsson 2020, Gawel 2017) wait for the approval batch.

\subsection{Capture proof, docking and magnetic grasping}
\label{sec:lit_capture}

A join needs three things the leave does not: a mechanism that catches, a proof that it has caught, and a switch
of control at the proof. The literature on in-flight docking and aerial grasping supplies precedents for each.

\paragraph{Proof of a catch.} A grasp can be proven by a sensor in the gripper or from the vehicles' motion. Shankar
et al.\ \cite{shankar2020inflight} confirm a bar grasp with a magnetic sensor before the giving vehicle releases
[Loianno 2018, pending the approval batch, for a sensed grasp of a static object]. ModQuad
\cite{saldana2018modquad} instead detects the new structure from motion capture, once the docked modules move as
one body, and only then uploads a new force distribution to every module; LEGION \cite{sugihara2026legion} holds a
zero-clearance interlock by controlling the contact force and torque and treats the latched joint as proven when
the wrench settles. A proof from pose is the only kind available to this thesis, since nothing on the ring or the
magnet is instrumented: the contact and the lift that follows it are read from motion capture alone
(Section~\ref{sec:reconfig}), and Section~\ref{sec:sim_attach} tests that proof against the rig's own noise.

\paragraph{Docking in flight.} Modular vehicles dock to one another in the air, often through magnets, and the
controller switch at contact matters as much as the mechanism. ModQuad docks with passive permanent magnets that
pull in small misalignments while the other structure hovers; TRADY \cite{sugihara2023trady} assembles and
disassembles through a switchable magnet and pegs, and cut the altitude jump after assembly from about 0.39 to
0.065\,m by blending the thrust from the old controller to the new one; Zhang et al.\ \cite{zhang2026epm} dock
micro-vehicles with an electro-permanent magnet that holds without power and switches on a current pulse. Miyazaki
et al.\ \cite{miyazaki2018airborne} dock a small multirotor onto a bar that a larger one lowers on a winch and then
raises, clear of its downwash, and the LEGION units dock end to end in flight and then grasp and carry objects as
one articulated body, outdoors. Each docks to a rigid structure or to a single carrier; none joins a payload hung
from a team, where the newcomer's pull at the plate can tilt the payload and the team's tensions must be
redistributed. Two further precedents take or shape a load from the air without docking to a vehicle: Antal et
al.\ \cite{antal2026hook} pick a load from one moving platform and set it down on another with a hook and a robust
adaptive MPC, on hardware, and Xu et al.\ \cite{xu2026hitches} control the hitch that two cables form around a
payload in midair, in simulation; neither changes a carrying team.

\paragraph{Magnetic grasping.} A magnet holds best when pulled straight off its plate and loses hold as the pull
turns across the face, so that the angle of pull, not only its magnitude, sets whether a magnetic grasp survives
a transient [Gawel et al.\ 2017, pending the approval batch, for magnetic picking with a multirotor]. The
rig's electromagnets catch on contact only and release under a jerk rather than under a steady load
(Section~\ref{sec:overview}); every ring lost on the rig began with a magnet letting go
(Section~\ref{sec:sim_detach}). That is why the rods are planned steep, why an announced change is ramped, and why
the newcomer is folded in from a small tension (Section~\ref{sec:reconfig}): each keeps every pull close to its
plate's axis and small in its rate of change.
